% GENERATED FILE. Edit canonical lessons, then run npm run generate:books.
% canonical-source-hash: fnv1a64:aa18f2ad9f05e63c
% canonical-lessons: HI-C471-vinamra, HI-C471-badtamiz, HI-C471-lihaz, HI-C471-sistacar, HI-C471-cirhana

\chapter{Manners}
\label{ch:hi-manners}
\hlchaptermodality{471}

\begin{chapteropening}
\textbf{By the end of this chapter:} \emph{I can say whether someone is polite or rude.}

Complete the last lesson of chapter 471: I can say whether someone is polite or rude.
\end{chapteropening}

\section[\texorpdfstring{\hi{विनम्र} (vinamra)}{विनम्र (vinamra)}]{\hi{विनम्र} (vinamra) — polite; humble}

\label{lesson:HI-C471-vinamra}



\begin{quote}
Before the new one: say the Hindi for fearless, then the Hindi for a risk.
\end{quote}

\subsection*{You'll want to know: \hi{विनम्र}}
\textbf{\hi{विनम्र}} — \emph{vinamra} — \textquotedblleft{}polite; humble\textquotedblright{}.

\begin{quote}
\hi{मीरा} \hi{बड़ों} \hi{से} \hi{हमेशा} \hi{विनम्र} \hi{होकर} \hi{बात} \hi{करती} \hi{है।}
\end{quote}

\begin{quote}
\emph{Mīrā baṛoṁ se hameśā vinamra hokar bāt kartī hai.}
\end{quote}

\begin{quote}
Mira always speaks politely to her elders.
\end{quote}

\begin{cousinweb}
\textbf{\hi{विनम्र}} (\emph{vinamra}) is a tatsama, borrowed from Sanskrit \emph{vinamra}, \textquotedblleft{}bent down, submissive, humble, modest\textquotedblright{} (Monier-Williams): someone who bows a little.
\end{cousinweb}

\begin{grammarlens}[title={What you've built}]
One more word for describing what you lived and felt.
\end{grammarlens}

\subsection*{Your turn}
\begin{itemize}
\raggedright
  \item \emph{Say it:} \emph{vinamra}
  \item \emph{Say it:} \emph{vinamra}, once more
  \item \emph{Say it:} \emph{Mīrā baṛoṁ se hameśā vinamra hokar bāt kartī hai.} — Mira always speaks politely to her elders.
  \item \emph{Recall:} say the Hindi for fearless, then the Hindi for a risk, then say \emph{vinamra} again
\end{itemize}

\begin{tcolorbox}[breakable,colback=teal!4,colframe=teal!35!black,title={Before you move on}]
What does \hi{विनम्र} mean? (\textquotedblleft{}Polite; humble\textquotedblright{}.) Say it once more.
\end{tcolorbox}
\section[\texorpdfstring{\hi{बदतमीज़} (badtamīz)}{बदतमीज़ (badtamīz)}]{\hi{बदतमीज़} (badtamīz) — rude; ill-mannered}

\label{lesson:HI-C471-badtamiz}



\begin{quote}
Before the new one: say the Hindi for a risk, then the Hindi for polite.
\end{quote}

\subsection*{You'll want to know: \hi{बदतमीज़}}
\textbf{\hi{बदतमीज़}} — \emph{badtamīz} — \textquotedblleft{}rude; ill-mannered\textquotedblright{}.

\begin{quote}
\hi{वह} \hi{लड़का} \hi{बहुत} \hi{बदतमीज़} \hi{है}, \hi{बड़ों} \hi{का} \hi{आदर} \hi{नहीं} \hi{करता।}
\end{quote}

\begin{quote}
\emph{vah laṛkā bahut badtamīz hai, baṛoṁ kā ādar nahīṃ kartā.}
\end{quote}

\begin{quote}
That boy is very rude; he doesn't respect his elders.
\end{quote}

\begin{cousinweb}
\textbf{\hi{बदतमीज़}} (\emph{badtamīz}) is usually traced to Persian \emph{bad}, \textquotedblleft{}bad\textquotedblright{}, joined to Arabic \emph{tamīz}, \textquotedblleft{}discernment, good manners\textquotedblright{}, which came through Persian: having bad manners.
\end{cousinweb}

\begin{grammarlens}[title={What you've built}]
One more word for describing what you lived and felt.
\end{grammarlens}

\subsection*{Your turn}
\begin{itemize}
\raggedright
  \item \emph{Say it:} \emph{badtamīz}
  \item \emph{Say it:} \emph{badtamīz}, once more
  \item \emph{Say it:} \emph{vah laṛkā bahut badtamīz hai, baṛoṁ kā ādar nahīṃ kartā.} — That boy is very rude; he doesn't respect his elders.
  \item \emph{Recall:} say the Hindi for a risk, then the Hindi for polite, then say \emph{badtamīz} again
\end{itemize}

\begin{tcolorbox}[breakable,colback=teal!4,colframe=teal!35!black,title={Before you move on}]
What does \hi{बदतमीज़} mean? (\textquotedblleft{}Rude; ill-mannered\textquotedblright{}.) Say it once more.
\end{tcolorbox}
\section[\texorpdfstring{\hi{लिहाज़} (lihāz)}{लिहाज़ (lihāz)}]{\hi{लिहाज़} (lihāz) — consideration; regard}

\label{lesson:HI-C471-lihaz}



\begin{quote}
Before the new one: say the Hindi for polite, then the Hindi for rude.
\end{quote}

\subsection*{You'll want to know: \hi{लिहाज़}}
\textbf{\hi{लिहाज़}} — \emph{lihāz} — \textquotedblleft{}consideration; regard\textquotedblright{}.

\begin{quote}
\hi{बूढ़े} \hi{लोगों} \hi{का} \hi{लिहाज़} \hi{करना} \hi{चाहिए।}
\end{quote}

\begin{quote}
\emph{būṛhe logoṁ kā lihāz karnā cāhiye.}
\end{quote}

\begin{quote}
One should show consideration for old people.
\end{quote}

Masculine.

\begin{cousinweb}
\textbf{\hi{लिहाज़}} (\emph{lihāz}) is usually traced to Arabic \emph{liḥāẓ}, \textquotedblleft{}a look, regard\textquotedblright{}, which reached Hindi through Persian: having regard for someone.
\end{cousinweb}

\begin{grammarlens}[title={What you've built}]
One more word for describing what you lived and felt.
\end{grammarlens}

\subsection*{Your turn}
\begin{itemize}
\raggedright
  \item \emph{Say it:} \emph{lihāz}
  \item \emph{Say it:} \emph{lihāz}, once more
  \item \emph{Say it:} \emph{būṛhe logoṁ kā lihāz karnā cāhiye.} — One should show consideration for old people.
  \item \emph{Recall:} say the Hindi for polite, then the Hindi for rude, then say \emph{lihāz} again
\end{itemize}

\begin{tcolorbox}[breakable,colback=teal!4,colframe=teal!35!black,title={Before you move on}]
What does \hi{लिहाज़} mean? (\textquotedblleft{}Consideration; regard\textquotedblright{}.) Say it once more.
\end{tcolorbox}
\section[\texorpdfstring{\hi{शिष्टाचार} (śiṣṭācār)}{शिष्टाचार (śiṣṭācār)}]{\hi{शिष्टाचार} (śiṣṭācār) — etiquette; correct manners}

\label{lesson:HI-C471-sistacar}



\begin{quote}
Before the new one: say the Hindi for rude, then the Hindi for consideration.
\end{quote}

\subsection*{You'll want to know: \hi{शिष्टाचार}}
\textbf{\hi{शिष्टाचार}} — \emph{śiṣṭācār} — \textquotedblleft{}etiquette; correct manners\textquotedblright{}.

\begin{quote}
\hi{मेहमान} \hi{के} \hi{आने} \hi{पर} \hi{खड़े} \hi{होना} \hi{शिष्टाचार} \hi{है।}
\end{quote}

\begin{quote}
\emph{mehmān ke āne par khaṛe honā śiṣṭācār hai.}
\end{quote}

\begin{quote}
Standing up when a guest arrives is good manners.
\end{quote}

Masculine.

\begin{cousinweb}
\textbf{\hi{शिष्टाचार}} (\emph{śiṣṭācār}) is a tatsama, borrowed from Sanskrit \emph{śiṣṭācāra}, \textquotedblleft{}the conduct of the learned, good manners\textquotedblright{} (Monier-Williams), built on \emph{śiṣṭa}, \textquotedblleft{}well taught\textquotedblright{}, and \emph{ācāra}, \textquotedblleft{}conduct\textquotedblright{}.
\end{cousinweb}

\begin{grammarlens}[title={What you've built}]
One more word for describing what you lived and felt.
\end{grammarlens}

\subsection*{Your turn}
\begin{itemize}
\raggedright
  \item \emph{Say it:} \emph{śiṣṭācār}
  \item \emph{Say it:} \emph{śiṣṭācār}, once more
  \item \emph{Say it:} \emph{mehmān ke āne par khaṛe honā śiṣṭācār hai.} — Standing up when a guest arrives is good manners.
  \item \emph{Recall:} say the Hindi for rude, then the Hindi for consideration, then say \emph{śiṣṭācār} again
\end{itemize}

\begin{tcolorbox}[breakable,colback=teal!4,colframe=teal!35!black,title={Before you move on}]
What does \hi{शिष्टाचार} mean? (\textquotedblleft{}Etiquette; correct manners\textquotedblright{}.) Say it once more.
\end{tcolorbox}
\section[\texorpdfstring{\hi{चिढ़ाना} (ciṛhānā)}{चिढ़ाना (ciṛhānā)}]{\hi{चिढ़ाना} (ciṛhānā) — to tease}

\label{lesson:HI-C471-cirhana}



\begin{quote}
Before the new one: say the Hindi for consideration, then the Hindi for etiquette.
\end{quote}

\subsection*{You'll want to know: \hi{चिढ़ाना}}
\textbf{\hi{चिढ़ाना}} — \emph{ciṛhānā} — \textquotedblleft{}to tease\textquotedblright{}.

\begin{quote}
\hi{बड़ा} \hi{भाई} \hi{मीरा} \hi{को} \hi{उसके} \hi{नाम} \hi{पर} \hi{चिढ़ाता} \hi{है।}
\end{quote}

\begin{quote}
\emph{baṛā bhāī Mīrā ko uske nām par ciṛhātā hai.}
\end{quote}

\begin{quote}
Mira's big brother teases her about her name.
\end{quote}

Past, with \hi{ने} (he or she …): \hi{उसने} \hi{बहन} \hi{को} \hi{चिढ़ाया}, \emph{usne bahan ko ciṛhāyā}.

\begin{cousinweb}
\textbf{\hi{चिढ़ाना}} (\emph{ciṛhānā}) is the causative partner of \emph{ciṛhnā}, \textquotedblleft{}to be provoked\textquotedblright{}, which CDIAL lists under \emph{*ciḍ} (CDIAL 4794): to make someone cross on purpose, usually in fun.
\end{cousinweb}

\begin{grammarlens}[title={What you've built}]
One more word for describing what you lived and felt.
\end{grammarlens}

\subsection*{Your turn}
\begin{itemize}
\raggedright
  \item \emph{Say it:} \emph{ciṛhānā}
  \item \emph{Say it:} \emph{ciṛhānā}, once more
  \item \emph{Say it:} \emph{baṛā bhāī Mīrā ko uske nām par ciṛhātā hai.} — Mira's big brother teases her about her name.
  \item \emph{Recall:} say the Hindi for consideration, then the Hindi for etiquette, then say \emph{ciṛhānā} again
\end{itemize}

\begin{tcolorbox}[breakable,colback=teal!4,colframe=teal!35!black,title={Before you move on}]
What does \hi{चिढ़ाना} mean? (\textquotedblleft{}To tease\textquotedblright{}.) Say it once more.
\end{tcolorbox}
